%Version 3.1 December 2024
% See section 11 of the User Manual for version history
%
%%%%%%%%%%%%%%%%%%%%%%%%%%%%%%%%%%%%%%%%%%%%%%%%%%%%%%%%%%%%%%%%%%%%%%
%%                                                                 %%
%% Please do not use \input{...} to include other tex files.       %%
%% Submit your LaTeX manuscript as one .tex document.              %%
%%                                                                 %%
%% All additional figures and files should be attached             %%
%% separately and not embedded in the \TeX\ document itself.       %%
%%                                                                 %%
%%%%%%%%%%%%%%%%%%%%%%%%%%%%%%%%%%%%%%%%%%%%%%%%%%%%%%%%%%%%%%%%%%%%%

%%\documentclass[referee,sn-basic]{sn-jnl}% referee option is meant for double line spacing

%%=======================================================%%
%% to print line numbers in the margin use lineno option %%
%%=======================================================%%

%%\documentclass[lineno,pdflatex,sn-basic]{sn-jnl}% Basic Springer Nature Reference Style/Chemistry Reference Style

%%=========================================================================================%%
%% the documentclass is set to pdflatex as default. You can delete it if not appropriate.  %%
%%=========================================================================================%%

%%\documentclass[sn-basic]{sn-jnl}% Basic Springer Nature Reference Style/Chemistry Reference Style

%%Note: the following reference styles support Namedate and Numbered referencing. By default the style follows the most common style. To switch between the options you can add or remove “Numbered” in the optional parenthesis. 
%%The option is available for: sn-basic.bst, sn-chicago.bst%  
 
%%\documentclass[pdflatex,sn-nature]{sn-jnl}% Style for submissions to Nature Portfolio journals
%%\documentclass[pdflatex,sn-basic]{sn-jnl}% Basic Springer Nature Reference Style/Chemistry Reference Style

\documentclass[pdflatex,sn-mathphys-num,iicol]{sn-jnl}% Math and Physical Sciences Numbered Reference Style

%%\documentclass[pdflatex,sn-mathphys-ay]{sn-jnl}% Math and Physical Sciences Author Year Reference Style
%%\documentclass[pdflatex,sn-aps]{sn-jnl}% American Physical Society (APS) Reference Style
%%\documentclass[pdflatex,sn-vancouver-num]{sn-jnl}% Vancouver Numbered Reference Style
%%\documentclass[pdflatex,sn-vancouver-ay]{sn-jnl}% Vancouver Author Year Reference Style
%%\documentclass[pdflatex,sn-apa]{sn-jnl}% APA Reference Style
%%\documentclass[pdflatex,sn-chicago]{sn-jnl}% Chicago-based Humanities Reference Style

%%%% Standard Packages
%%<additional latex packages if required can be included here>

\usepackage{graphicx}%
\usepackage{multirow}%
\usepackage{amsmath,amssymb,amsfonts}%
\usepackage{amsthm}%
\usepackage{mathrsfs}%
\usepackage[title]{appendix}%
\usepackage{xcolor}%
\usepackage{textcomp}%
\usepackage{manyfoot}%
\usepackage{booktabs}%
\usepackage{algorithm}%
\usepackage{algorithmicx}%
\usepackage{algpseudocode}%
\usepackage{listings}%
\usepackage{comment}%
%%%%

\usepackage{draftwatermark}

\newcommand{\av}[1]{\langle #1 \rangle}
\newcommand{\blue}[1]{\textcolor{blue}{#1}}
\newcommand{\red}[1]{\textcolor{red}{#1}}

%%%%%=============================================================================%%%%
%%%%  Remarks: This template is provided to aid authors with the preparation
%%%%  of original research articles intended for submission to journals published 
%%%%  by Springer Nature. The guidance has been prepared in partnership with 
%%%%  production teams to conform to Springer Nature technical requirements. 
%%%%  Editorial and presentation requirements differ among journal portfolios and 
%%%%  research disciplines. You may find sections in this template are irrelevant 
%%%%  to your work and are empowered to omit any such section if allowed by the 
%%%%  journal you intend to submit to. The submission guidelines and policies 
%%%%  of the journal take precedence. A detailed User Manual is available in the 
%%%%  template package for technical guidance.
%%%%%=============================================================================%%%%

%% as per the requirement new theorem styles can be included as shown below
\theoremstyle{thmstyleone}%
\newtheorem{theorem}{Theorem}%  meant for continuous numbers
%%\newtheorem{theorem}{Theorem}[section]% meant for sectionwise numbers
%% optional argument [theorem] produces theorem numbering sequence instead of independent numbers for Proposition
\newtheorem{proposition}[theorem]{Proposition}% 
%%\newtheorem{proposition}{Proposition}% to get separate numbers for theorem and proposition etc.

\theoremstyle{thmstyletwo}%
\newtheorem{example}{Example}%
\newtheorem{remark}{Remark}%

\theoremstyle{thmstylethree}%
\newtheorem{definition}{Definition}%

\raggedbottom
%%\unnumbered% uncomment this for unnumbered level heads

\begin{document}

\title[Experimental Investigation of Classical Billiards Dynamics Using an Autonomous Robot and Image Analysis]{Experimental Investigation of Classical Billiards Dynamics Using an Autonomous Robot and Image Analysis}

%%=============================================================%%
%% GivenName	-> \fnm{Joergen W.}
%% Particle	-> \spfx{van der} -> surname prefix
%% FamilyName	-> \sur{Ploeg}
%% Suffix	-> \sfx{IV}
%% \author*[1,2]{\fnm{Joergen W.} \spfx{van der} \sur{Ploeg} 
%%  \sfx{IV}}\email{iauthor@gmail.com}
%%=============================================================%%

\author[1]{\fnm{Ruan V. A.} \sur{Quirino}}\email{ruan.quirino@ufrpe.br}

\author[1]{\fnm{Francisco C. B.} \sur{Leal}}
\email{fbonleal@gmail.com}

\author[1]{\fnm{Antonio R. de C.} \sur{Romaguera}}
\email{antonio.romaguera@ufrpe.br}

\author[1]{\fnm{Anderson L. R.} \sur{Barbosa}}
\email{anderson.barbosa@ufrpe.br}

\author*[1]{\fnm{T. Araújo} \sur{Lima}}\email{tiago.araujol@ufrpe.br}
\equalcont{These authors contributed equally to this work.}

\affil[1]{\orgdiv{Departamento de Física}, \orgname{Universidade Federal Rural de Pernambuco}, \orgaddress{\city{Recife}, \postcode{52171-900}, \state{Pernambuco}, \country{Brazil}}}

%%==================================%%
%% Sample for unstructured abstract %%
%%==================================%%

\abstract{

\ \ \ \ The construction of a versatile experimental system, integrating technologies such as Arduino, modeling, 3D printing, and image processing, proved effective in investigating the dynamics of active particles. Our system enabled the exploration of billiards dynamics, a traditional nonlinear mechanical system in which a particle is free to move within a closed domain, experiencing elastic collisions on the boundary. In recent years, these systems have been experimentally studied in optical and robotic systems. In general, quantitative analyses are based on trajectories in the phase space~(positions and velocities). To ensure precision, efficiency, and reliability in these measures, we choose specific components as motors, distance sensors, voltage controllers, batteries, velocity controllers, and cameras. The experimental results showed good agreement with theoretical predictions, validating the methodology employed. Furthermore, the developed system has potential for expansion to study new geometric scenarios and the effects of external perturbations, consolidating its utility as an experimental tool for analyzing complex dynamic systems.

}

\keywords{Billiards, Classical Dynamics, Robot, Lyapunov Exponent}

%%\pacs[JEL Classification]{D8, H51}

%%\pacs[MSC Classification]{35A01, 65L10, 65L12, 65L20, 65L70}

\maketitle

\section{Introduction}
\label{intro}

\ \ \ \ This work describes a macroscopic experimental setup developed to observe the dynamical behavior of a mathematical billiard system. A digital camera recorded the interaction between a robot and the billiards walls, where specular reflections occur. Dynamical systems are a pillar of classical physics, these formulations describe the temporal evolution of a set of state variables~\cite{ott:2002}. For continuous time, the evolution is described by differential equations. These equations could be solved analytically in some cases. In general, they are nonlinear and coupled, making them untractable with analytical tools~\cite{str:2018}. In the 19th century, H. Poincaré introduced geometrical methods to study dynamical systems. The focus of his new approach is to examine the structural properties of solutions rather than exact solutions. The methods consist, in sum, of identifying fixed points, limit cycles, and stable or unstable varieties in the phase space. H. Poincaré also demonstrated that small perturbations can lead to significant divergences in trajectories~\cite {poi:1899}. Today, this phenomenon is known as \textit{Deterministic Chaos.} The phase space contains the entire set of possible system configurations. It can exhibit three main behaviors: periodic, quasi-periodic, and chaotic. The periodic behavior means that the state variables $x_i (t + 2\pi/\omega_i) = x_i(t)$ and the set $\{ \omega_i\}$ is comensurable. For quasi-periodic case $\{ \omega_i\}$ is uncomensurable. Chaotic systems are characterized by sensitivity to initial conditions. The exponential separation of trajectories with nearby initial conditions leads to the quantity called \textit{Lyapunov Exponent~($\lambda$)}. For $\lambda>0$, the system is chaotic, indicating that the trajectories are apart from each other as $\exp(\lambda t)$. Here, $\lambda$ is the largest Lyapunov exponent~\cite{ben:1976}. A double pendulum~\cite{stg:1992}, and the Lorentz system are examples of chaotic systems~\cite{lor:1963}.

Discrete-time systems work by iteractions where the present state depends on the previous state from a deterministic function, defining a map. These systems do not need discretization algorithms as in continuos systems. There are several examples of discrete-times chaotic systems, including one-dimensional cases as the logistic map~\cite{pnp:2006}. Two-dimensional systems most commonly cite are the Henon map~\cite{bmc:1991}, the standard map, and the billiards tables, central theme of this paper. The standard map, also known as Chirikov-Taylor map~\cite{chi:1979}, works as a simples model to explore the transition from regularity to chaos in conservative dynamical systems. It models the behavior of a kicked rotor, a system where a rod is free to rotate around an axis between periodic impulses. For small values of a control parameter, the trajectories of the phase space are confined. In this regime, trajectories are periodic or quasi-periodic. The parameter increasing leads to the emergence of the called \textit{Kolmogorov-Arnold-Moser~(KAM) Islands}, or just \textit{stability islands} and chaotic regions~\cite{chi:1979,ott:2002}. Thus, in the same phase space coexist initial conditions with regular dynamics and chaotic ones, forming so a \textit{Mixed Phase Space}~\cite{ree:1999,ara:2015,ara:2024,qui:2025}.

Billiards are toy models in the study of nonlinear dynamical systems. In billiard tables, a particle is free to move within a closed domain between elastic collisions at the boundary. The system is described by two position variables and two velocities, forming a four-dimensional phase space $(x,y,v_x,v_y)$. This dynamics is more relevant in the collision moments. Birkhoff coordinates are more convenient to describe the dynamics in a billiard~\cite{bir:2020}. This pair of phase space coordinates $(\ell,\sin\phi)$ is computed just at the collision points along the billiard border~(Fig.~\ref{fig:coord_birkhoff}). The particle position along the border is parametrized by $\ell$, while $\phi$ is the incidence angle. So, $\sin\phi$ represents the tangent velocity at the border. The border form is fundamental for the billiard dynamics. The behavior could vary from regular to chaotic just with tiny perturbations on the border geometry~\cite{leo:2019}. In a circular billiard, there is no sensitivity to initial conditions. Thus, tiny perturbations on the particle position and/or velocity do not result in significant changes in the particle trajectory~\cite{dre:1998}. In this case, the tangent velocity has the role of angular momentum, which is conserved in this dynamics as the mechanical energy. So, the phase space described by the Birkhoff coordinates consists of just a countable number of points or a single line along the constant-tangent-velocity curve for periodic or quasi-periodic dynamics, respectively. These properties serve as the circular billiard basis for testing and calibrating the experimental setup described in this work. This paper is organized as follows. Section~II presents the experimental setup and data acquisition. In Sec.~III, we describe the experimental procedure and trajectory reconstruction. The results and their analysis are discussed in Secs.~IV and V. Finally, Sec.~VI summarizes the main conclusions and perspectives.

\begin{figure}[!htbp]
\centering
\includegraphics[width=0.95\columnwidth]{Fig1.pdf}
\caption{\label{fig:coord_birkhoff} Illustration of Birkhoff coordinates: $\ell$ is the arc length along the edge, and $\sin \phi$ represents the tangential velocity at a collision.}
\end{figure}

\section{\label{sec:met} Experimental Setup}

\subsection{\label{subsec:rob_desc}Robot Description}

\ \ \ \ The robot developed is an autonomous mobile platform designed to simulate the behavior of a point particle in classical billiards systems, as shown in Fig.~\ref{fig:robot}. Built with reproducible components, the robot utilizes an Arduino Nano microcontroller~(ATmega328P) \cite{NANO:2024} as its central processing unit, which is responsible for controlling the motors, reading the sensors, and implementing the reflection laws. The movement is generated by two continuous rotation DC servomotors~(FS5103R)~\cite{FS5103R}, powered by a regulated voltage of $6$~V through an LM2596 DC-DC converter~\cite{LM2596}, ensuring stable speed and adequate torque. The system power is provided by four 18650 batteries connected in series~\cite{santanaimport_datasheet}, with a total voltage of approximately $15$~V, regulated to $9$~V for Arduino and $6$~V for motors by LM2596 modules~\cite{lm2596s_datasheet}, ensuring safe and efficient operation. Table~\ref{tab:comp_simple} lists all electronic components used in the robot's construction, detailing their main functions.

\begin{figure}[!htbp]
\centering
\includegraphics[width=0.75\columnwidth]{Fig2.pdf}
\caption{\label{fig:robot} Real image of the robot manufactured using 3D printing, with the electronic components installed.}
\end{figure}

\begin{table*}[!htbp]

\caption{\label{tab:comp_simple} Main functions of the electronic components used in the robot.}

\begin{tabular}{| p{0.22\textwidth} | p{0.73\textwidth} |}

\hline

\textbf{Component} & \textbf{Main Function} \\

\hline
\hline

Arduino Nano microcontroller (ATmega328P) & Acts as the central processing unit, controlling motors, reading sensors, and implementing reflection laws. \\

\hline

DC servo-motors (FS5103R) & Generates the continuous rotation movement and propulsion for the mobile platform. It is responsible for the rotation of the platform. \\

\hline

LM2596 DC-DC converter & Regulates and stabilizes the voltage to 6~V to provide adequate torque and constant motor speed. \\

\hline

18650 batteries & Provides the total power to the system (approximately 15~V when connected in series). \\

\hline

LM2596 modules & Regulates the battery voltage to specific levels: 9~V for the Arduino and 6~V for the motors. \\

\hline

LM393 optical encoder-based & Functions as a speed sensor for precise velocity control and trajectory correction. \\

\hline

VL53L0X distance sensor & Detects the billiard's edges at distances under 15~cm using Time-of-Flight (ToF) technology. \\

\hline

\end{tabular}

\end{table*}

For precise speed control and trajectory correction, the robot is equipped with LM393 optical encoder-based speed sensors~\cite{lm393_speed_sensor}. A central element of the design is the wheel with an integrated encoder disk, 3D-printed from polylactide~(PLA) and featuring 90~perforations evenly distributed along its edge. This configuration allows for an angular resolution of $2^\circ$ per detection event, as each transition between the perforation and the solid segment generates a pulse in the sensor. The schematic in Fig.~\ref{fig:encoder} illustrates how the system works. This integrated design eliminates additional mechanical couplings, reduces alignment errors, and significantly improves motion control accuracy. In addition, the robot uses three VL53L0X distance sensors~\cite{VL53L0X}, based on Time-of-Flight~(ToF) technology, arranged at fixed angles. Two lateral at $45^\circ$ and one central at $0^\circ$, for precise detection of the billiard's edges at distances of less than $15$~cm. These sensors offer high resolution, low latency, and good immunity to environmental interference, such as lighting variations.

\begin{figure}[!htbp]
\centering
\includegraphics[width=0.75\columnwidth]{Fig3.pdf}
\caption{\label{fig:encoder} Operating diagram of an optical encoder: disc with markings, light emitter and receiver.}
\end{figure}

The robot's mechanical structure is 3D printed from PLA using a 3D printer, ensuring high dimensional accuracy and reproducibility. The main chassis is almost circular, with a diameter of $14$~cm, and is optimized for dynamic stability and low weight. Fig.~\ref{fig:3Dmodels}(a) shows the wheels designed to be directly coupled to the engine axis. The main chassis, Fig.~\ref{fig:3Dmodels}(b), was designed to berth components such as the Arduino controller, batteries, and engines. Guaranteeing a low center of mass. To close the chassis's top, a lid was designed (Fig.~\ref{fig:3Dmodels}(c)). It protects the circuits and receives the LEDs in the robot's geometrical center. Figs.~\ref{fig:3Dmodels}(d) and (e) show the support for sensors, encoders, and distance, respectively. Support and stabilizing parts are shown in Figs.~\ref{fig:3Dmodels}(f), (g), and (h). To avoid mechanical vibrations and to generate an alignment between motors, axes, and wheels, we employ a support column shown in Fig.~\ref{fig:3Dmodels}(i). The 3D models of all parts are designed in FreeCAD~\cite{freecad} and are available in \texttt{.STL}~\cite{qui:2026rep}. The structural components are attached with $3$~mm screws threaded directly into the preheated PLA, a process that creates robust internal threads and ensures mechanical rigidity without the need for nuts or adhesives. The Arduino microcontroller and auxiliary circuits are mounted on a perforated universal copper board, facilitating secure electrical connections, access for maintenance, and quick module replacement if necessary. The robot has a total mass of $780$~g and is programmed to move in a straight line at a velocity $|v| \simeq 20$~cm/s until it detects an edge of the billiards at a distance of less than $15$~cm, at which point it stops the movement, calculates the angle of reflection and performs a rotation to restart the trajectory with an exit angle equal to the angle of incidence. Thus, simulating the behavior of a point particle in a classical billiard system.

\begin{figure*}[!htpb]
\centering
\includegraphics[width=0.95\textwidth]{Fig4.pdf}
\caption{\label{fig:3Dmodels} Three-dimensional models developed in FreeCAD~\cite{freecad}, and rendered using Blender. (a) Wires. (b) Main chassi. (c) Structural cover. (d) Couppling element to receive encoders. (e) Support to sensors of distance. (f) Support collumn whose connect the main chassi with top parts. (g) Stabilizar support. (h) Structure for switches. (i) Support for engines.}
\end{figure*}

\subsection{\label{sec:bil_set} Billiards Settings}

\ \ \ \ The experimental billiards were constructed to achieve a high degree of geometric precision, reproducibility, and mechanical stability, essential for the controlled study of dynamic systems. Although different geometries can be investigated with the same robot, such as the lemon billiards~\cite{qui:2025}, or the Bunimovich stadium~\cite{vas:2022}, this section describes the general method of assembly and contour calibration, taking the circular billiards as an example. A fundamental configuration for system validation and characterization of the robot's reflexive behavior. The billiards' outline is formed by flexible MDF~(medium-density fiberboard) boards, with a thickness of $6$~mm, a height of $20$~cm, and adjustable length, fixed to the laboratory floor using 3D-printed joints and locks~(model available at~\cite{qui:2026rep}). These structural elements guarantee that the billiards' shape remains unchanged during experiments, minimizing vibrations and misalignments. Assembly begins with the precise positioning of reference points on the floor, defined using geometric drawing tools~(such as hand compasses or laser guides), to maintain symmetry and alignment with the video acquisition system.

In the case of circular billiards, the radius $R$ is pre-calculated based on the desired total perimeter $L_0$, kept constant between different configurations to allow direct comparisons between geometries. The radius is given by $R = L_0 / 2\pi$, ensuring that the perimeter of the circle is comparable to that of other non-circular configurations. With the center of the circle fixed, a compass is used to trace the outline on the floor. This trace serves as a visual guide for the exact positioning of the boards, which are curved and fixed to faithfully follow the drawn line. Alignment accuracy is critical to minimize systematic errors in the robot's trajectories, especially during reflections. To achieve this, the acquisition camera~(GoPro HERO9) is positioned vertically above the center of the billiards, to obtain an orthogonal view of the surface. Small tilts are digitally corrected during image analysis, but the initial parallel alignment between the camera, the floor, and the contour significantly reduces the need for post-processing corrections.

The final geometric validation is performed through image analysis: a static image of the billiards is captured and processed to determine the actual contour shape with millimeter precision. This step allows the identification of local deviations~(such as unwanted curvatures or gaps in the joints) and the calculation of effective parameters that are used in comparisons with theoretical models. Among these parameters, the effective radius $R_\text{eff}$ stands out, defined as the adjusted value of the contour's radius of curvature, obtained from the analysis of the real trajectories traveled by the robot during the experiments. Fig.~\ref{fig:Exp_bil} illustrates a real image of the circular billiards, with the robot's trajectory superimposed. The orange dots indicate the detected collision positions, whose fit to a circle allows the determination of the effective boundary corresponding to $R_\text{eff}$.

\begin{figure}[!htbp]
\centering
\includegraphics[width=0.9\columnwidth]{Fig5.pdf}
\caption{\label{fig:Exp_bil} Trajectory data~(in black) and collisions~(in orange), superimposed on the photograph of the experiment corresponding to a circular billiard.}
\end{figure}

To determine the velocity, we used data extracted from images captured by the camera. Initially, the robot's trajectory was identified, and the camera's capture rate was considered. Then, the trajectory dimensions, originally expressed in pixels, were normalized to metric units. The velocity was calculated as the ratio of the spatial displacement to the time interval between consecutive points. In Fig.~\ref{fig:velocity}, the evolution of the speed over the first minute of the experiment is observed. The orange curve represents the velocity filtered with a low-pass filter, highlighting smoother fluctuations and facilitating the identification of abrupt velocity drops, which correspond to collision moments. Collision detection is directly related to its velocity. When the robot reaches the edge, the system interrupts the movement, calculates the necessary parameters, and performs a reflection, a rotation around its own axis. During this rotation, the centralized LED tends to remain fixed in place, showing no apparent displacement. This behavior is identified as a collision, detected by the abrupt drop in velocity, and the corresponding position is recorded for analysis. In Fig.~\ref{fig:Exp_bil}, it is possible to see the collision points, demonstrating the effectiveness of the procedure, by observing that all collisions are in fact on the edges of the circular billiard.

\begin{figure}[!htbp]
\centering
\includegraphics[width=0.9\columnwidth]{Fig6.pdf}
\caption{Instantaneous velocity in the first minute. The orange curve shows the low-pass filtered version.}
%Bottom: Robot trajectory (black dots, one per frame, totaling ~54,000). Detected collisions are highlighted by larger orange dots with a black outline.}
\label{fig:velocity}
\end{figure}

\subsection{\label{sec:exp_proc} Experimental Procedure}

\ \ \ \ The experimental data are acquired using a GoPro HERO9 Black camera positioned vertically above the center of the billiard, ensuring an orthogonal view of the surface. The camera was configured to record in Full HD resolution~($1920 \times 1080$~pixels) at a frame rate of 30~fps, using the Narrow field of view mode. This setting was chosen to minimize fisheye geometric distortions, preserving the fidelity of the billiard table's shape and allowing for precise spatial measurements. Each experimental run lasted approximately 30~minutes, resulting in videos with approximately 54,000~frames per experiment. All experiments were conducted in darkness to maximize the contrast between the white LED and the environment, while the white LED centered on the robot acts as a point source of highly contrasting light. This experimental condition eliminates external light interference and unwanted reflections, facilitating the automatic detection of the robot's position with high precision. The use of active lighting only on the robot functions as a passive optical marker, ideal for intensity-based tracking algorithms.

The recorded videos were processed offline using an automated pipeline based on FFmpeg and Python/MATLAB scripts. The routines used in each step are available in~\cite{qui:2026rep}. First, the files are decomposed into sequences of individual frames, converted to grayscale, and organized chronologically. Next, a mask is applied to limit the analysis to the physical area of the billiard. The robot's position in each frame is determined by a peak light intensity detection algorithm, whose output is metrically calibrated based on a previously registered reference grid. The final data consists of time-stamped $(x,y)$ position time series, allowing for complete trajectory reconstruction and precise identification of events such as collisions.

The robot is programmed to simulate the behavior of an ideal particle in a classic billiard system. During movement, it travels in a straight line until it detects the edge of the billiard table at a distance of less than 15~cm, using its ToF sensors~(VL53L0X). Upon detecting a collision, the robot immediately stops moving and calculates the angle of reflection $\theta_r$ equal to the angle of incidence $\theta_i$, based on the relative orientation between its trajectory and the unit vector normal to the wall. These angles ate obtained by a geometrical procedure ilustrated in Fig.~\ref{fig:ref_rule}. Vector $\vec{S}_1$ lies on the robot movement direction, while $\vec{S}_2$, and $\vec{S}_3$ determine the boundary in the collisio region. The tangent vector off border in the collision region is given by
%
\begin{equation}
\label{eq:tanvec}
\hat{e}_T = \frac{\vec{S}_3 - \vec{S}_2}{\left | \vec{S}_3 - \vec{S}_2 \right |}.
\end{equation}
%
The normal direction is given by a $\pi/2$ rotation over the tangent vector, $\hat{e}_N = \hat{R}_{\pi/2} \cdot \hat{e}_T$. In this new basis~(nb), vector $\vec{S}_1$ is expressed as
%
\begin{equation}
\label{eq:newbasis}
\vec{S}_1^\text{(nb)} = \left ( \vec{S}_1 \cdot \hat{e}_N \right ) \hat{e}_N + \left ( \vec{S}_1 \cdot \hat{e}_T \right ) \hat{e}_T.
\end{equation}
%
In this coordinate system, the incident angle is just $\theta_i = \pi - \varphi$ as in Fig.~\ref{fig:ref_rule}. Where
%
\begin{equation}
\label{eq:ref_angle}
\varphi = \cos^{-1} \left ( \frac{\left | \vec{S}_1^\text{(nb)} \right |}{\vec{S}_1 \cdot \hat{e}_N} \right  ).
\end{equation}
%
It then rotates around its own axis until it reaches the new angle of displacement and restarts linear motion. To ensure straight trajectories and constant velocity~(20~cm/s), the control uses feedback from the optical encoders~(LM393), individually adjusting the rotation of the wheels~(see Fig.~\ref{fig:encoder}). This well-defined dynamic, combined with precise trajectory tracking, allows for the acquisition of robust experimental data for the analysis of dynamic systems.

\begin{figure}[!htbp]
\centering
\includegraphics[width=0.85\columnwidth]{Fig7.pdf}
\caption{\label{fig:ref_rule} Scheme of robot's sensors and geometrical procedures to calculate $\theta_i$.}
\end{figure}

\section{Results}

\subsection{\label{subsec:cali} Calibration}

\ \ \ \ With the robot's construction completed, the billiard boundary structure finalized, and the data acquisition capability fully established, the next step is to evaluate the robot's efficiency in performing the experiments. To this end, three fundamental aspects were analyzed: Boundary detection, straight-line displacement, and specular reflections. The calibration is performed on a circular billiard because its dynamics conserve the angle of incidence. Or in other words, the angular momentum is conserved. Thus, choosing a circular billiard to calibrate the setup not only facilitates the analysis of reflections, but also provides a reliable basis for validating the system before transitioning to experiments with chaos. First of all, we need to test the system's ability to accurately identify collision events, ensuring that no relevant interaction was missed during the experiment. Fig.~\ref{fig:Exp_bil} shows data points at the collisions events at a circular border.

The second step is to test the robot's ability to maintain straight trajectories, a criterion of high importance in the construction and calibration of the system. Between each collision, during the robot's free flight, the system is programmed to maintain a constant wheel rotation, as detailed in previous sections. To verify whether the robot's displacement is aligned with the ideal, we adopted a trajectory-based method. Initially, a numerical straight line is drawn between the points of consecutive collisions, assuming that the ideal displacement would be perfectly aligned with this line. The quantitative analysis of this alignment is performed using the coefficient of determination $(R^2)$, which assesses the proportion of variability in the robot's position explained by this ideal trajectory. $R^2$ is defined as:
%
\begin{equation}
\label{eq:R2}
R^2 = 1 - \frac{\sum_{i=1}^n \left ( \vec{r}_i - \vec{r}_i^\text{(ideal)} \right )^2}{\sum_{i=1}^n \left ( \vec{r}_i - \vec{r}_i^\text{(mean)} \right )^2}.
\end{equation}
%
Where $\vec{r}_i$ is the actual position of the robot, $\vec{r}_i^\text{(ideal)}$ is the position predicted by the ideal line, and $\vec{r}_i^\text{(mean)}$ is the mean of the observed positions. High $R^2$ values indicate that the robot moved consistently and aligned with the ideal trajectory, while low values suggest deviations. This method provides a robust metric for evaluating the robot's movement accuracy and adjusting the control system, ensuring the reliability of rectilinear motion between collisions. As illustrated in Fig.~\ref{fig:R2}, $ \av{R^2} = 0.999 \pm 0.002$ for an experimental run demonstrates that the robot's trajectory is aligned with the ideal, with small deviations attributable to experimental noise. This analysis confirms the effectiveness of the implemented control and allows a quantitative evaluation of the quality of the flights.

\begin{figure}[!htbp]
\centering
\includegraphics[width=0.975\columnwidth]{Fig8.pdf}
\caption{Representation of 250 rectilinear trajectories with the coefficient of determination $R^2$ calculated for each one. Values are displayed in the interval $[0.975,\, 1]$, allowing a detailed analysis of the small variations in $R^2$.}
\label{fig:R2}
\end{figure}

Finally, it must to confirm that the robot's reflection angle corresponds to the angle of incidence, to characterize reflections close to ideal. Analysis of reflections is fundamental for verifying that the robot maintains specular behavior throughout the experiment, which is essential to guarantee the validity of the results. With the ability to consistently detect the edges of the billiard table and synchronize the wheels' speeds, it becomes necessary to evaluate the accuracy of the reflections. This analysis justifies the choice of a circular billiard for the calibration and initialization of the setup. In the upper panel of Fig.~\ref{fig:RDelta_Phi}, the ratio $R_{\phi,i}$ between the largest and smallest angles between incidence and reflection is observed over about 300~collisions. The average $\av{R_\phi} = 0.98 \pm 0.01$ indicates that the robot's reflections are close to a specular behavior, with no significant deviations. In the lower pannel of Fig.~\ref{fig:RDelta_Phi}, the difference between the angles of consecutive collisions, $\delta \phi_i = |\phi_i - \phi_{i-1}|$ is shown. This analysis highlights the constancy of reflection, showing that, in a circular billiard, the theoretical variation of the angle should be zero. The experimental data corroborate this expectation, presenting an average $\av{\delta_\phi} = 0.8 \pm 0.6^\circ$. These results are fundamental to validating the robot in more complex experiments, since gross reflection errors could compromise the analysis of chaotic dynamics. Therefore, using a circular billiard in the calibration step allows us to ensure the absence of chaos and to evaluate the robot's performance in detail. In a circular billiard, it is expected that the motion will be regular with constant angular momentum. In the experimental phase space, the trajectories follow the numerical predictions, with a noise of approximately 0.1 in angular momentum. Fig.~\ref{fig:phase_space} illustrates this regime.

\begin{figure}[!htbp]
\centering
\includegraphics[width=0.975\columnwidth]{Fig9.pdf}
\caption{Upper panel: $R_{\phi}$ as a function of collision $i$. Values close to 1 indicate reliable specular reflections. Lower panel: Difference between the current reflection angle and the previous reflection angle in degrees $(\delta \phi = |\phi_{i} - \phi_{i-1} |)$, evaluating the consistency of the system during the experiment.}
\label{fig:RDelta_Phi}
\end{figure}

\begin{figure}[!htbp]
\centering
\includegraphics[width=0.975\columnwidth]{Fig10.pdf}
\caption{\label{fig:phase_space} Phase space of collision points in the robot trajectory on a circular billiard, as described in Fig.~\ref{fig:Exp_bil}.}
\end{figure}

\subsection{\label{subsec:lyap} Lyapunov Exponent}

\ \ \ \ The Lyapunov exponent $(\lambda)$ is a quantitative measure to access the sensitivity to initial conditions in dynamical systems. It reflects the behavior of nearby trajectories in phase space and how their separation diverges over time. Being particularly relevant in the current study, the calculation of $\lambda$ was defined by the expression:
%
\begin{equation}
\label{eq:lyapdef}
\lambda = \lim_{T \rightarrow \infty}\frac{1}{T}\sum_{i=0}^{N-1} \frac{\left | \vec{d}_{i+1} \right |}{\left | \vec{d}_i \right |},
\end{equation}
%
where $\left | \vec{d}_i \right |$ represents the norm of the separation vector between two trajectories at collision $i$, $T = \sum_{i=1}^{N} \tau_i$, and $\tau_i$ is the continuous time between consecutive collisions. The calculation begins with the definition of the initial separation vector $\vec{d}_0 = (d_\ell, d_s)$ whose components are chosen randomly and then normalized to ensure that $\left | \vec{d}_0 \right | = 1$. \textcolor{red}{(LEMBRAR DE DEFINIR COORD. DE BIRKORV. E COLOCAR S = SIN(PHI) LÁ EM CIMA. DEFINIR "d" COMO PERTURBAÇÃO. E DIZER QUE "J" É JACOBIANO.)} This initialized vector will undergo successive transformations that reflect the evolution of the dynamic system. At each collision, $\vec{d}$ is updated using evolution matrices and is normalized after each step to avoid divergence. The dynamics of the separations are modeled by two matrices. The first one, $\hat{J}_0$, describes the free flight between two collisions in a time interval $\tau$, being represented by:
%
\begin{equation}
\label{eq:J0}
\hat{J}_0 = 
\begin{bmatrix}
1 & \tau \\
0 & 1
\end{bmatrix}.
\end{equation}
%
This matrix represents uniform rectilinear motion, in which $d_\ell$ receives an increment of $\tau d_s$, while $d_s$ remains unchanged. Next, the matrix $\hat{J}_\text{c}$, responsible for capturing the effects of specular reflections, is applied to the transformed vector. The matrix $\hat{J}_\text{c}$ is expressed as:
%
\begin{equation}
\label{eq:Jc}
\hat{J}_\text{c} = -\begin{bmatrix}
1 & 0 \\
\frac{2 \kappa}{\cos \phi)} & 1
\end{bmatrix},
\end{equation}
%
where $\kappa$ represents the local curvature and $\phi$ the angle of reflection. The total transformation over a collision is obtained by combining these two matrices. This matrix describes the complete evolution of the vector $\vec{d}$ in a collision, accounting for the accumulated effects of free flight and border interaction. Over $N$ collisions, $\vec{d}$ evolves successively, with matrices adapted at each step as a function of the parameters $\kappa_i$, $\phi_i$, and $\tau_i$ obtained experimentally. At each collision, the norm of $\vec{d}$ is calculated and is renormalized to avoid excessive growth during subsequent calculations~\cite{ben:1978,pik:2016,vas:2022,qui:2025}. Finally, the measure of Lyapunov exponent is given by:
%
\begin{equation}
\label{eq:lyapexp}
\lambda_\text{exp} = \frac{1}{T} \sum_{i=1}^N \ln \left | \vec{d}_i \right |.
\end{equation}
%
This calculation provides an accurate characterization of the system's sensitivity to initial conditions. For the trajectory shown in Figs.~\ref{fig:Exp_bil} and~\ref{fig:phase_space}, $\lambda_\text{exp} \simeq 0.002$. From an average of 10 measurements, $\av{\lambda_\text{exp}} = 0.012 \pm 0.008$. Each experimental run contains around 500~collisions. For comparison, in a numerical calculus using the same parameters, the result is $\av{\lambda_\text{num}} = 0.0076 \pm 0.0001$. The experimental result shows an excellent agreement with the theoretical prediction. Similar setups and procedures were used in chaotic~\cite{vas:2022} and mixed~\cite{qui:2025} billiards successfully.

\section{Conclusions}
\label{sec:conc}

\ \ \ \ The construction of a versatile experimental system, involving technologies such as Arduino, modeling, 3D printing, and image processing, proved effective for exploring the dynamics of active particles in complex systems. This setup demonstrated the feasibility of integrating robotics and computational analysis to investigate dynamic regimes in a controlled manner. The approach employed stood out for its adaptability, allowing the analysis of different geometric and dynamic scenarios. Furthermore, the developed system has expansion potential and could be adapted to investigate the effects of external perturbations and explore new geometric scenarios.

The development of the experimental system encompassed the conception and assembly of an Arduino-based robot with modules such as distance sensors and servomotors to the implementation of C/C++ for Arduino control and Python for data analysis. A camera captured videos of the robot's dynamics, which served as a data source for quantitative analysis. We used image processing techniques in MATLAB and Python to extract relevant information, such positions and velocities. 3D modeling enabled the precise design of the components required for the experiment. The experiments were performed in a domain with a fixed perimeter of approximately 5~m. To determine the experimental Lyapunov exponent $(\lambda_\text{exp})$, we used the tangent space-based method~\cite{ben:1978,pik:2016,vas:2022,qui:2025} and compared the experimental results with numerical predictions. The results demonstrated a satisfactory agreement between the experimental data and the theoretical approaches.

Our experimental setup successfully reproduced a highly idealized model by programming the robot's responses to adhere to its validity conditions. Robot-environment interaction provides an excellent platform for experimental studies in physical systems~\cite{sch:1998,dau:2019,leo:2020,li:2022,li:2023,gol:2024,par:2024}, including additional aspects of chaos and dynamical systems, such as self-propelled confined particles~\cite{dam:2022} and self-excluding active particles~\cite{alb:2024}. Furthermore, the study of robot swarms has been receiving considerable attention from the scientific community~\cite{gar:2018}. Robots can mimic not only the behavior of a confined billiard particle but also the movement of a school of fish in an aquarium~\cite{lea:2024,rom:2024,deSou:2024}. In this way, it is possible to study the collective dynamics of certain animal groups using robots~\cite{lei:2020}. Thus, this setup reinforce the relevance of experimental platforms integrating robotics and computing, offering tools for advanced studies in nonlinear physics and dynamical systems.

\bmhead{Acknowledgements}

\red{A VER.}

\section*{Declarations}

\red{CONFERIR, PEGUEI DE ARTIGO ANTIGO.}

\begin{itemize}

\item This work was supported in part by the following Brazilian agencies: Conselho Nacional de Desenvolvimento Cient\'ifico e Tecnol\'ogico (CNPq), under Grant Numbers 404057/2021-7 and 309457/2021 (ALRB).
\item The authors declare no conflict of interest.
\item Ethics approval and consent to participate: Not applicable.
\item Consent for publication: Ok.
\item Data availability: The data generated for the current study are available in~\cite{qui:2026rep}.
\item Materials availability: Not applicable.
\item Code availability: Codes are available in~\cite{qui:2026rep}.

\end{itemize}

\bibliography{refs_2026_InstrumentationBilliards}% common bib file
%% if required, the content of .bbl file can be included here once bbl is generated
%%\input sn-article.bbl

\end{document}